\documentclass[10pt,a4paper]{amsart}
\usepackage[margin=19mm]{geometry}
\usepackage{amsmath,amssymb,mathtools}
\usepackage[scr=rsfs,cal=euler]{mathalfa}
\usepackage{xfrac}
\usepackage[unicode,hidelinks,bookmarks=false]{hyperref}
\newcommand{\ie}{i.e.\ }
\newcommand{\cf}{cf.\ }
\newcommand{\resp}{respectively\ }
\newcommand{\Subsection}{\subsection}
\providecommand{\nobreakdash}{\nobreak\hbox{-}}
\setlength{\emergencystretch}{2em}
\multlinegap0pt
\usepackage{aligned-overset}
\DeclareMathOperator{\Exp}{Exp}
\DeclareMathOperator{\supp}{supp}
%\DeclareMathOperator{\dt}{dt}

\renewcommand*\r{\mathbb{R}}
\newcommand*\n{\mathbb{N}}
\newcommand*\z{\mathbb{Z}}
\newcommand*\one{\mathbf{1}}
\newcommand*\p{\mathbb{P}}
\newcommand*\all{\forall}
\newcommand*\eps{\varepsilon}
\newcommand*\tensor{\otimes}
\newcommand*\indexx[2]{n\left(#1,#2\right)}
\newcommand*\nonnegative{\left[0,\infty\right)}
\newcommand*\cal[1]{\mathcal{#1}}
\newcommand*\union{\cup}
\newcommand*\abs[1]{\left|#1\right|}

\newcommand*\intersect{\cap}


\def\hypergeometric#1#2#3#4{_{2}F_{1}\left(#1,#2;#3;#4\right)}
%%%%%%%%%%%%%%%%%%%%%%%%%%%%%%%%%%%%%%%%%%%%%%%%%%%%%%%%
%~ This fixes spacing issues around \left( and \right) and other similar symbols
\let\originalleft\left 
\let\originalright\right
\renewcommand{\left}{\mathopen{}\mathclose\bgroup\originalleft}
\renewcommand{\right}{\aftergroup\egroup\originalright}
%%%%
\renewcommand*{\to}{\mathchoice{\longrightarrow}{\rightarrow}{\rightarrow}{\rightarrow}}
%%%%
\let\oldin\in
\renewcommand{\in}{\mathchoice{\oldin}{\oldin}{\,\oldin\,}{\,\oldin\,}}

%%%%%%%%%%%%%%%%%%%%%%%%%%%%%%%%%%%%%%%%%%%%%%%%%%%%%%%%%%%%
\graphicspath{{./figures/}}

\newcommand*{\mk}{\mkern -1mu}
\newcommand*{\Mk}{\mkern -2mu}
\newcommand*{\mK}{\mkern 1mu}
\newcommand*{\MK}{\mkern 2mu}

\hypersetup{urlcolor=purple, linkcolor=blue, citecolor=red}

\newcommand*{\relabel}{\renewcommand{\labelenumi}{(\theenumi)}}
\newcommand*{\romanenumi}{\renewcommand*{\theenumi}{\roman{enumi}}\relabel}
\newcommand*{\Romanenumi}{\renewcommand*{\theenumi}{\Roman{enumi}}\relabel}
\newcommand*{\alphenumi}{\renewcommand*{\theenumi}{\alph{enumi}}\relabel}
\newcommand*{\Alphenumi}{\renewcommand*{\theenumi}{\Alph{enumi}}\relabel}
\let\oldtilde\tilde
\renewcommand*{\tilde}[1]{\mathchoice{\widetilde{#1}}{\widetilde{#1}}{\oldtilde{#1}}{\oldtilde{#1}}}
\let\oldforall\forall
\renewcommand*{\forall}{\mathrel{\oldforall}}

\newtheorem{theo}{Theorem}[section]
\newtheorem{prop}[theo]{Proposition}
\newtheorem{coro}[theo]{Corollary}
\newtheorem{rem}[theo]{Remark}
\newtheorem{rema}[theo]{Remark}
\theoremstyle{definition}\newtheorem{defi}[theo]{Definition}
\newenvironment{proofc}[1][]{}{}
\title[Discrete distributions are iid-distance laws]{The distance problem via subadditivity: original Theorem1.1}
\author{Renan Gross}
\date{Source: Annales Henri Lebesgue8(2025),1023--1035; DOI10.5802/ahl.254}
\begin{document}\maketitle
\noindent\textit{Editorial scope.} The counted unit is original Theorem1.1 with its complete author proof and the complete source-local selection-space mechanism it uses. Discreteness, nonnegativity and support at zero remain exact. The illustrative Figure2.1 is omitted with a bounded-source fidelity receipt because the cached package contains no figure asset; its single in-text reference becomes a link to the published article and no mathematical content is removed. Printed slips, including the conditioning event at205, the missing subscripts at265/269 and the run-together display at294, are recorded without generated corrections.
\section{Introduction}
\subsection{The problem}
The \emph{distance problem} asks the following. Given a distribution $\theta$ on $\nonnegative$, is it possible to find a complete separable metric space $(S,d)$ and a Borel probability measure $\mu$ on $S$, so that when $X,Y\sim\mu$ are independent, their distance satisfies $d(X,Y)\sim\theta$? If so, we say that $(S,d,\mu)$
\emph{achieves} $\theta$, and that $\theta$ is \emph{feasible}. This very natural and quite inviting problem was first proposed by Aldous, Blanc and Curien in~\cite{aldous_blanc_curien_distance_problem_2024}, where they also give some statistical motivation.

\subsection{Examples}
\begin{enumerate}
\setcounter{enumi}{2}
\item \emph{Bernoulli distribution:}\label{enu:bernoulli_construction} let $0< p_{0} < 1$. In order to achieve the Bernoulli distribution $\theta=p_{0}\delta_{0}+(1-p_{0})\delta_{1}$, let $m\geq2$ be an integer and let $\alpha\in[0,1]$ to be chosen later. Let $S=\{ 0,\ldots,m-1\} $, equipped with the discrete metric
\[
d\left(x,y\right)=
\begin{cases}
0 & x=y\\
1 & x\neq y.
\end{cases}
\]
Let $\mu$ be the distribution of the random variable on $S$ whose value is $0$ with probability $\alpha$, and uniform among $\{ 1,\ldots,m-1\} $ with probability $1-\alpha$. Then
\begin{equation}\label{eq:bernoulli_probability}
\p\left[d\left(X,Y\right)=0\right]=\alpha^{2}+\left(\frac{1-\alpha}{m-1}\right)^{2}.
\end{equation}
Choosing $m$ so that the quadratic equation $p_{0}=\alpha^{2}+(\frac{1-\alpha}{m-1})^{2}$ has a solution and solving for $\alpha$ gives us the desired probability space.

\end{enumerate}
\subsection{Results}
\begin{theo}\label{thm:discrete_distributions}
Every nonnegative discrete distribution $\theta$ supported on $0$ is feasible. \hyperlink{target:proof_of_discrete}{({proof $\nearrow$})}%revoir la destination ? 
\end{theo}

\section{Proof of Theorem~\ref{thm:discrete_distributions}}

\hypertarget{target:proof_of_discrete}Our proof uses the following construction, which, given several feasible distributions, returns a mixture of these distributions. The resultant distance can be thought of as being computed by a ``selection'' mechanism: we linearly order the distributions, and iteratively flip a coin for each one to see whether or not it is the one that we sample from; if none of the distributions are chosen, the default $0$ is returned.

\subsection{Selection space}

Denote the integers by $\z$ and the positive integers by $\n=\{ 1,2,\ldots\} $. We say that a triplet $(S,d,\mu)$ is a measured metric space if $(S,d)$ is a metric space and $\mu$ is a Borel probability measure on $S$ with respect to $d$.

Let $I\subseteq\z$, and for each $n\in I$, let $(S_{n},d_{n},\mu_{n})$ be a complete countable measured metric space with more than $1$ point. Assume that for every $n\in I\intersect\n$ there exists a special $s_{n}\in S_{n}$ such that
\begin{equation}\label{eq:prob_sums_are_finite}
\sum_{n\in I\,\intersect\,\n}1-\mu_{n}\left(\left\{ s_{n}\right\} \right)<\infty.
\end{equation}
We now define a new space $(S,d,\mu)$ out of these spaces; later, in Proposition~\ref{prop:selection_space_properties}, we give conditions under which it is a complete separable measured metric space.
\begin{enumerate}
\item \emph{The space}: we say that a vector $x\in\prod_{n\in I}S_{n}$ is \emph{eventually constant} if there exists an $N>0$ such that $x_{n}=s_{n}$ for all $n\geq N$ (we still index the vector $x$ by the set $I$). Note that when $\sup I<\infty$, every vector is eventually constant. We set
\[
S=\left\{ x\in\prod_{n\in I}S_{n}\,\middle| \,x\text{ is eventually constant}\right\}.
\]
\item \emph{The function }$d$: for $x\neq y\in S$, let $\indexx xy=\sup\{ n\mid x_{n}\neq y_{n}\} $. This is always finite since $x$ and $y$ are eventually constant. We then define
\[
d\left(x,y\right)=
\begin{cases}
0 & x=y\\
d_{\indexx xy}\left(x_{\indexx xy},y_{\indexx xy}\right) & x\neq y.
\end{cases}
\]
\item \emph{The measure}: let $\tilde{\mu}=\tensor_{n\in I}\mu_{n}$ be the product measure on $\prod_{n\in I}S_{n}$. By~\eqref{eq:prob_sums_are_finite} the sum $\sum_{n\in I\intersect\n}1-\mu_{n}(\{ s_{n}\} )$ is finite, and the Borel--Cantelli lemma implies that $X\sim\tilde{\mu}$ is eventually constant with probability $1$ under $\tilde{\mu}$. The measure $\tilde{\mu}$ can therefore be restricted to a measure $\mu$ on $S$.
\end{enumerate}
We call the tuple $(S,d,\mu)$ the \emph{selection space} of $\{ (S_{n},d_{n},\mu_{n})\} _{n\in I}$. Let us find the distribution of $d(X,Y)$ when $X,Y\sim\mu$ are independent. The probability of obtaining $0$ is
\begin{equation}\label{eq:prob_of_0}
\p\left[d\left(X,Y\right)=0\right]=\prod_{n\in I}\p\left[X_{n}=Y_{n}\right]=\prod_{n\in I}\p\left[d_{n}\left(X_{n},Y_{n}\right)=0\right],
\end{equation}
while for $r>0$,
\begin{multline}\label{eq:prob_of_positive}
\p\left[d\left(X,Y\right)=r\right] \\
\begin{aligned}
& =\sum_{n\in I}\p\left[n\left(X,Y\right)=n\right]\cdot\p\left[d_{n}\left(X_{n},Y_{n}\right)=r\,\middle|\,\indexx XY=n\right]\\
& =\sum_{n\in I}\left(\prod_{k\in I,\,k\,>\,n}\p\left[X_{k}=Y_{k}\right]\right)\p\left[X_{n}\neq Y_{n}\right]\cdot\p\left[d_{n}\left(X_{n},Y_{n}\right)=r\,\middle|\, X_{n}\neq Y_{n}\right] \\
& =\sum_{n\in I}\left(\prod_{k\in I,\,k\,>\,n}\p\left[d\left(X_{k},Y_{k}\right)=0\right]\right)\p\left[d_{n}\left(X_{n},Y_{n}\right)=r\right],
\end{aligned}
\end{multline}
and we have indeed obtained a mixture of positive elements of the distributions achieved by $(S_{n},d_{n},\mu_{n})$.
\begin{rem}
As a space of vectors indexed by $I\subseteq\z$, the selection space can also be interpreted as the steps of a bi-directional walk, where at time-step $n$, the walker may go in the ``direction'' $x_{n}\in S_{n}$. When $\sup I$ is finite and we view time as starting at $\sup I$ and directed down towards $\inf I$, the index $\indexx xy$ is the first time that two walkers $x$ and $y$ split up when starting at a common origin. In this case the elements of $S$ can be seen as the space of ends on some rooted tree, and the function $d$ is a generalization of the metric $e^{-h(x,y)}$ on the space of ends, where $h(x,y)$ is the depth of the last common vertex of the rays $x$ and $y$. See \href{https://ahl.centre-mersenne.org/item/AHL_2025__8__1023_0/}{Figure 2.1}. When $\sup I=\infty$, however, there is no such common root.

\end{rem}

\begin{prop}\label{prop:selection_space_properties}
Let $I\subseteq\z$ and $\{ (S_{n},d_{n},\mu_{n})\} _{n\in I}$ be as above. Assume that the following two properties hold.
\begin{enumerate}
\item\label{prop2.2.1} For every $m<n\in I$, every $x_{1}\neq x_{2}\in S_{m}$ and every $y_{1}\neq y_{2}\in S_{n}$,
\begin{equation}\label{eq:metrics_are_subadditive_assumption}
d_{m}\left(x_{1},x_{2}\right)\leq2d_{n}\left(y_{1},y_{2}\right).
\end{equation}
\item\label{prop2.2.2} We have
\begin{equation}\label{eq:assumption_separability}
\text{either }
\inf I>-\infty\quad
\text{or}\quad
\lim_{n\to-\infty}\,\sup_{z,w\in S_{n}}d_{n}\left(z,w\right)=0.
\end{equation}
\end{enumerate}
Then the selection space $(S,d,\mu)$ is a complete separable measured metric space.
\end{prop}

\begin{proofc}
\begin{proof}[{Metric}]
It is clear that $d(x,y)=0$ if and only if $x=y$, so we only need to verify the triangle inequality. Let $x,y,z\in S$ be distinct, and denote $b=\indexx xy$.
\begin{enumerate}
\item If either $\indexx yz>b$ or $\indexx xz>b$, then since $x$ and $y$ agree on indices greater than $b$, we must have $\indexx xz=\indexx yz>b$, and also $d(x,z)=d(y,z)$. Then by~\eqref{eq:metrics_are_subadditive_assumption},
\[
d\left(x,y\right)=d_{b}\left(x_{b},y_{b}\right)\leq2d_{\indexx xz}\left(x_{\indexx xz},z_{\indexx xz}\right)=d\left(x,z\right)+d\left(y,z\right).
\]
\item If both $n(y,z)=b$ and $n(x,z)=b$, then the triangle inequality is satisfied because $d_{b}$ is a metric.
\item Otherwise, we have $n(y,z),n(x,z)\leq b$, and at least one of $n(y,z)$ or $n(x,z)$ is strictly smaller than $b.$ But it is impossible for both to be strictly smaller, since that would mean that $x_{b}=z_{b}=y_{b}$, contradicting the fact that $x_{b}\neq y_{b}$. Supposing w.l.o.g. that $\indexx xz=b$, we then trivially have
\[
d\left(x,y\right)=d_{b}\left(x_{b},y_{b}\right)=d_{b}\left(x_{b},z_{b}\right)=d\left(x,z\right)\leq d\left(x,z\right)+d\left(y,z\right).
\]
\end{enumerate}
\let\qed\relax
\end{proof}
\begin{proof}[{Completeness}]
Let $N\in I$, and let $a\neq b\in S_{N}$. Then by~\eqref{eq:metrics_are_subadditive_assumption}, for every $n>N$ and every $w,z\in S_{n}$, $d_{n}(w,z)\geq\frac{1}{2}d_{N}(a,b)$. Thus, for every $N\in I$,
\begin{equation}
\inf_{n>N}\inf_{w,z\in S_{n}}d_{n}\left(w,z\right)>0.\label{eq:inf_on_tail}
\end{equation}
Let $(x^{(k)})_{k\in\n}$ be a Cauchy sequence in $S$, and let $L=\inf I$. When $L=-\infty$, \eqref{eq:inf_on_tail} implies that if $d(x^{(k_{1})},x^{(k_{2})})<\eps$ for some $k_{1}$ and all $k_{2}\geq k_{1}$ and small enough $\eps$, then $x_{n}^{(k_{2})}=x_{n}^{(k_{1})}$ for all $n>N$ for some $N\in I$, with $N\to-\infty$ as $\eps\to0$. Thus $x^{(k)}$ converges to an element in $S$. When $L>-\infty$, \eqref{eq:inf_on_tail} implies that for small enough $\eps$, $x_{n}^{(k_{2})}=x_{n}^{(k_{1})}$ for all $n>L$, while the first elements $(x_{L}^{(k)})_{k\in\n}$ form a Cauchy sequence in $S_{L}$. Since $S_{L}$ is itself complete, we again get that $x^{(k)}$ converges to an element in~$S$.\let\qed\relax
\end{proof}
\begin{proof}[{Separability}] 
If $\inf I>-\infty$, then $S$ is a countable union of countable sets, and so is countable itself. Otherwise, by~\eqref{eq:assumption_separability}, $\lim_{n\to-\infty}\sup_{z,w\in S_{n}}d(z,w)=0.$ For each negative $n\in I$, let $a^{(n)}\in S_{n}$, and set
\[
A_{n}=\left\{ x\in S\,\middle|\, x_{i}=a^{\left(i\right)}\,\all i\leq n\right\}.
\]
Since all vectors in $S$ are eventually constant, each $A_{n}$ is a countable union of countable sets and is therefore countable, and so the union $\union_{n\in I}A$ is also countable. It is also dense in $S$: for any $x\in S$ and any $n\in I$, the set $A_{n}$ contains an element $y$ such that $y_{k}=x_{k}$ for all $k>n$, and so $n(x,yt)\leq n$; denseness follows since $\sup_{z,w\in S_{n}}d_{n}(z,w)\to0$ as $n\to-\infty$.\let\qed\relax
\end{proof}

\begin{proof}[{Borel measure}]
Under the metric $d$, an open ball around $x\in S$ of radius $r$ is given by
\[
B\left(x,r\right)=\bigcup_{N\in I}\quad\bigcup_{s\in S_{N}\mid d_{N}\left(x_{N},s\right)<r,s\neq x_{N}}\left\{ y\in S\,\middle|\, y_{N}=s,y_{n}=x_{n}\text{ for all }n>N\right\} \union\left\{ x\right\}.
\]
Each set of the form $\{ y\in S\mid y_{N}=s,y_{n}=x_{n}\text{ for all } n>N\} $ is measurable under the product measure $\tilde{\mu}$, and therefore so is the countable union of such sets. The restriction measure $\mu$ is therefore defined on the $\sigma$-algebra generated by the open balls in $S$. Since $(S,d)$ is separable, $\mu$ is Borel.
\end{proof}\let\qed\relax
\end{proofc}

\begin{prop}\label{prop:positive_interval_is_feasible}
Let $a>0$, let $p_{0}\in(0,1)$
and let $\lambda$ be a discrete probability distribution with $\supp(\lambda)\subseteq[a,2a].
$
Then the distribution $\theta=p_{0}\delta_{0}+(1-p_{0})\lambda$ is feasible.
\end{prop}

\begin{proof}
Let $a_{0}=0$, and let $(a_{n})_{n\in I}$ be the positive atoms of $\lambda$ arranged in some arbitrary order; here, $I=\n$ if there is an infinite number of atoms and $I=[N]$ for some $N$ otherwise. For every $n\in I$, denote $p_{n}=\theta(\{ a_{n}\})$, and define $q_{n}$ by
\[
q_{n}=\frac{p_{n}}{p_{0}+p_{1}+\ldots+p_{n}}.
\]
For $n\in I$, let $(S_{n},d_{n},\mu_{n})$ be the metric space corresponding to the Bernoulli measure $\theta_{n}=(1-q_{n})\delta_{0}+q_{n}\delta_{a_{n}}$, obtained by multiplying the metric described in the examples section by $a_{n}$. Then the assumption in the construction of the selection space and the two conditions in Proposition~\ref{prop:selection_space_properties} hold for $\{ (S_{n},d_{n},\mu_{n})\} _{n\in I}$:
\begin{enumerate}
\item When constructing a Bernoulli random variable $p\delta_{0}+(1-p)\delta_{1}$with probability $p$ of obtaining $0$ as in Example~\ref{enu:bernoulli_construction}, the parameters $m$ and $\alpha$ may be chosen so that $\alpha\geq p$. Indeed, setting $f_{m}(x)=x^{2}+(\frac{1-x}{m-1})^{2}$, we have $p=f_{m}(\alpha)$. Since $\lim_{m\to\infty}f_{m}(p)=p^{2}<p$, there exists an $m$ large enough so that $f_{m}(p)<p$, and by continuity of $f_{m}$ and the fact that $f_{m}(1)=1$, there exists $\alpha\in[p,1]$ with $f_{m}(\alpha)=p$ as required. This shows that the probability to obtain $0$ under $\mu_{n}$ is bounded by
\begin{equation}\label{eq:bound_on_q}
\mu_{n}\left(\left\{ 0\right\} \right)\geq1-q_{n}.
\end{equation}
Thus
\begin{equation}\label{eq:sum_of_q_is_finite}
\sum_{n\in I}1-\mu_{n}\left(\left\{ 0\right\} \right)\leq\sum_{n\in I}q_{n}=\sum_{n\in I}\frac{p_{n}}{p_{0}+p_{1}+\ldots+p_{n}}\leq\sum_{n\in I}\frac{p_{n}}{p_{0}}=\frac{1-p_{0}}{p_{0}},
\end{equation}
and the sum in~\eqref{eq:prob_sums_are_finite} is finite.
\item Since $a_{n}\in[a,2a]$, the requirement~\eqref{eq:metrics_are_subadditive_assumption} is immediate: the left-hand side of~\eqref{eq:metrics_are_subadditive_assumption} is at most $2a$, while the right-hand side is at least $2a$.
\item $\inf I=1>-\infty$ as required in~\eqref{eq:assumption_separability}.
\end{enumerate}
Thus, the selection space $(S,d,\mu)$ constructed from $\{(S_{n},d_{n},\mu_{n})\} _{n\in I}$ is a complete separable measured metric space. Since the positive atoms of the $\mu_{n}$s are all distinct, equations~\eqref{eq:prob_of_0} and~\eqref{eq:prob_of_positive} imply that
\begin{equation}\label{eq:single_prob_zero}
\p\left[d\left(X,Y\right)=0\right]=\prod_{n\in I}\p\left[d_{n}\left(X_{n},Y_{n}\right)=0\right]=\prod_{n\in I}\left(1-q_{n}\right)=p_{0},
\end{equation}
and
\[
\p\left[d\left(X,Y\right)=a_{n}\right]=q_{n}\prod_{i\in I,\,i\,>\,n}\left(1-q_{i}\right)=p_{n}
\]
as needed.

We remark that by~\eqref{eq:bound_on_q} and~\eqref{eq:single_prob_zero},
\begin{equation}\label{eq:entire_vector_is_0_whp}
\p\left[X=\left(0,0,\ldots\right)\right]=\prod_{n\in I}\mu_{n}\left(\left\{ 0\right\} \right)\geq\prod_{n\in I}\left(1-q_{n}\right)=p_{0}.\qedhere
\end{equation}
\end{proof}

\begin{proof}[Proof of Theorem~\ref{thm:discrete_distributions}]\label{prfTheo1.1}
The measure $\theta$ can be written as
\[
\theta=p_{\infty}\delta_{0}+\sum_{n\in I}p_{n}\lambda_{n},
\]
where $I\subseteq\z$, and for every $n\in I$, $p_{n}>0$ and the measure $\lambda_{n}$ is discrete and supported on $(2^{n},2^{n+1}]$. Note that if $p_{\infty}=0$, then necessarily $\inf I=-\infty$. Define the numbers
\[
\beta_{n}=\frac{p_{n}}{1-\sum_{i\in I,\,i\,>\,n}p_{i}},
\]
and let $(S_{n},d_{n},\mu_{n})$ be the metric space defined in the proof of Proposition~\ref{prop:positive_interval_is_feasible} corresponding to the measure $\theta_{n}=(1-\beta_{n})\delta_{0}+\beta_{n}\lambda_{n}$. Note that each $S_{n}$ is countable. Then the assumption in the construction of the selection space and the two conditions in Proposition~\ref{prop:selection_space_properties} hold for $\{(S_{n},d_{n},\mu_{n})\} _{n\in I}$:
\begin{enumerate}
\item If $\sup I<\infty$ then of course the sum in~\eqref{eq:prob_sums_are_finite} is finite. Otherwise, by~\eqref{eq:entire_vector_is_0_whp}, if $X_{n}\sim\mu_{n}$, then
\[
\p\left[X_{n}=\left(0,0,\ldots\right)\right]\geq1-\beta_{n}.
\]
Thus
\begin{align*}
\sum_{n\in I\,\intersect\,\n} \p\left[X_{n}\neq\left(0,0,\ldots\right)\right]&\leq\sum_{n\in I\,\intersect\,\n}1-\left(1-\beta_{n}\right) =\sum_{n\in I\,\intersect\,\n}\beta_{n} \\ &\leq\sum_{n\in I\,\intersect\,\n}\frac{p_{n}}{1-\sum_{i\in I,\,i\,>\,n}p_{i}}\leq\frac{1}{p_{\infty}+\sum_{i\in I,\,i\,<\,1}p_{i}}.
\end{align*}
\item Since the range of $d_{n}$ is $(2^{n},2^{n+1}]$, the requirement~\eqref{eq:metrics_are_subadditive_assumption} is immediate; in fact, the left hand side of~\eqref{eq:metrics_are_subadditive_assumption} is bounded by just half of the right hand side.
\item $\lim_{n\to-\infty}\sup_{z,w\in S_{n}}d_{n}(z,w)=\lim_{n\to-\infty}2^{n+1}=0$ as required in~\eqref{eq:assumption_separability}.
\end{enumerate}
Thus, the selection space $(S,d,\mu)$ constructed from $\{(S_{n},d_{n},\mu_{n})\} _{n\in I}$ is a complete separable measured metric space. Since the supports of $\lambda_{n}$ are disjoint, equations~\eqref{eq:prob_of_0} and~\eqref{eq:prob_of_positive} imply that
\[
\p\left[d\left(X,Y\right)=0\right]=\prod_{n\in I}\p\left[X_{n}=Y_{n}\right]=\prod_{n\in I}\left(1-\beta_{n}\right)=p_{\infty},
\]
while the probability to sample from $\lambda_{n}$ is exactly $\beta_{n}\prod_{i\in I,\,i\,>\,n}(1-\beta_{i})=p_{n}$, as needed.
\end{proof}

\begin{thebibliography}{99}
\bibitem[ABC24]{aldous_blanc_curien_distance_problem_2024} D.J.Aldous, G.Blanc and N.Curien, ``The distance problem on measured metric spaces'',2024, https://arxiv.org/abs/2403.10926.
\end{thebibliography}
\end{document}
